% ---
% school: "college"
% level: "8eme"
% section: null
% chapter: "العمليات والحساب على الأعداد الكسرية"
% source: "Axiomaths"
% title: "أسئلة متعددة الخيارات"
% country: "Tunisie"
% year: 2026
% professor: "Axiomaths"
% difficulty: "Facile"
% points: null
% tags: ["qcm", "fractions", "operations"]
% ---
% Exercice 1/15 de la série 8ème « العمليات والحساب على الأعداد الكسرية » (statement + \axsol + solution). Compiler avec XeLaTeX + axiomath.sty (+ axiomaths-logo.png).
\documentclass[12pt,a4paper]{article}
\usepackage[level=8b,color=15803D]{axiomath}
\axsetup{title={العمليات والحساب على الأعداد الكسرية}, corrige}

\begin{document}

\begin{corr}[title={أسئلة متعددة الخيارات}, stars=1, source={Axiomaths}]
لكل سؤال إجابة صحيحة واحدة أو أكثر.
\begin{qcm}
\qitem[AB]{العدد $-\dfrac{3}{5}$ يساوي:}{$\dfrac{-3}{5}$}{$\dfrac{3}{-5}$}{$\dfrac{-3}{-5}$}{$-\left(-\dfrac{3}{5}\right)$}
\qitem[B]{قيمة المجموع $\dfrac{1}{2}+\dfrac{1}{3}$ هي:}{$\dfrac{2}{5}$}{$\dfrac{5}{6}$}{$\dfrac{1}{5}$}{$\dfrac{2}{6}$}
\qitem[B]{مقلوب العدد $-\dfrac{2}{7}$ هو:}{$\dfrac{2}{7}$}{$-\dfrac{7}{2}$}{$\dfrac{7}{2}$}{$-\dfrac{2}{7}$}
\qitem[ABD]{الأعداد التي تساوي $\dfrac{2}{3}$ هي:}{$\dfrac{14}{21}$}{$\dfrac{10}{15}$}{$\dfrac{9}{12}$}{$\dfrac{22}{33}$}
\qitem[A]{حاصل القسمة $\dfrac{3}{4}\div\dfrac{9}{8}$ يساوي:}{$\dfrac{2}{3}$}{$\dfrac{27}{32}$}{$\dfrac{3}{2}$}{$\dfrac{1}{6}$}
\qitem[AC]{الأعداد الأصغر من $-\dfrac{1}{2}$ هي:}{$-\dfrac{2}{3}$}{$-\dfrac{1}{3}$}{$-\dfrac{3}{4}$}{$\dfrac{1}{4}$}
\qitem[AC]{الجذاء $\dfrac{2}{3}\times\left(-\dfrac{9}{4}\right)$ يساوي:}{$-\dfrac{3}{2}$}{$\dfrac{3}{2}$}{$-\dfrac{18}{12}$}{$-\dfrac{8}{27}$}
\end{qcm}
\axsol
\axanswerkey{1/AB, 2/B, 3/B, 4/ABD, 5/A, 6/AC, 7/AC}
\begin{itemize}[itemsep=8pt]
  \item س1 : علامة الناقص يمكن أن تكون في البسط أو في المقام أو أمام الكسر : $\dfrac{-3}{5}=\dfrac{3}{-5}=-\dfrac{3}{5}$.
        أمّا $\dfrac{-3}{-5}=\dfrac{3}{5}$ و $-\left(-\dfrac{3}{5}\right)=\dfrac{3}{5}$ فهما موجبان.
  \item س2 : $\dfrac{1}{2}+\dfrac{1}{3}=\dfrac{3}{6}+\dfrac{2}{6}=\dfrac{5}{6}$. لا نجمع البسطين مع بعضهما ولا المقامين مع بعضهما.
  \item س3 : مقلوب $\dfrac{a}{b}$ هو $\dfrac{b}{a}$ وله نفس الإشارة : $-\dfrac{7}{2}$.
  \item س4 : $\dfrac{14}{21}=\dfrac{2\times7}{3\times7}$ و $\dfrac{10}{15}=\dfrac{2\times5}{3\times5}$ و $\dfrac{22}{33}=\dfrac{2\times11}{3\times11}$ ، أمّا $\dfrac{9}{12}=\dfrac{3}{4}$.
  \item س5 : $\dfrac{3}{4}\div\dfrac{9}{8}=\dfrac{3}{4}\times\dfrac{8}{9}=\dfrac{24}{36}=\dfrac{2}{3}$.
  \item س6 : $-\dfrac{2}{3}=-\dfrac{4}{6}$ و $-\dfrac{3}{4}=-\dfrac{9}{12}$ و $-\dfrac{1}{2}=-\dfrac{6}{12}$ ؛ و $-\dfrac{1}{3}=-\dfrac{2}{6}$ أكبر من $-\dfrac{1}{2}=-\dfrac{3}{6}$.
  \item س7 : $\dfrac{2}{3}\times\left(-\dfrac{9}{4}\right)=-\dfrac{18}{12}=-\dfrac{3}{2}$.
\end{itemize}
\end{corr}

\end{document}
